\section{2025-01-29}

\startframedtipp
{\bf Elektromagnetischer Schwingkreis}
\stopframedtipp

Vier aufgaben, drei davon müssen wir machen. Jede Aufgabe zu einem Rahmenthemer und einen Themenübergriff. Man kann kein Themengebiet raus lassen

Themen die sich anbieten (Auflistung!! aus den Tafelbildern):

Auswahl:

\startitemize[packed,m]
\item Überfliegen
\item Markieren, was kann ich lösen? Für {\bf jede} Teilaufgabe
\startitemize[packed]
\item Kann ich das sehr gut?
\item Kann ich das mit nach denken?
\item Kann ich das garnicht?
\stopitemize
\stopitemize

\startframedtipp
\startitemize[packed]
\item Elektromagnetischer Schwingkreis wird drann kommen
\item Irgendwas mit Doppelspalt
\startitemize[packed]
\item Gleichung herleiten
\item Aufbau
\stopitemize
\item Themen die noch noch nicht in normalen Klausuren dran kam
\item Aufbabe zum Thema bewerten wird von ihm hochgeladen
\item Grundsätzlich keine Formeln auswendig wissen
\item Herleitungen müssen ggf. nachgeschlagne werden. Ergebnisse von Herleitungen 
\item Geschwindigkeit geladenen Teilchen elektrische Felder wird drann kommen
\item Elektronenkanone
\item Wellenlaenge am gitter berechnen
\stopitemize
\stopframedtipp

\startframedtipp
\startitemize[packed,m]
\item Ein oder zwei Abende die Inhalte aus dem vorheriegen Unterricht angucken
\item
\stopitemize
\stopframedtipp

\startframedtipp

\startplaceformula[reference]
\startformula

\stopformula
\stopplaceformula
\stopframedtipp

\section{2025-01-29 - Modell Atomhülle}

Potenzialtopf (unendlich hoh, eindimensional)

\startframedtipp
Idee: Ein eindimensionaler Topf mit unendlich hohen Wänden, wird mit Atomen gefüllt.
\stopframedtipp

Energie $e^-$ in der Atomhülle ist in festen Niveaus, also quantisiert.

\startbuffer[modell_atomhuelle:2025:01:29:skizze]
\startMPpage

u := 1cm;

draw (-1.5u, 3u)--(-1.5u, 0)--(1.5u, 0)--(1.5u, 3u);

label.top("Wand", (-1.5u, 3u));
label.top("Wand", (1.5u, 3u));
label.top(btex $e^-$ etex, (0, 3u));
label.bot(btex $L$ etex, (0, 0));

label.lft(btex n etex, (-1.5u, 3u));
label.lft(btex 3 etex, (-1.5u, 2u));
label.lft(btex 2 etex, (-1.5u, 1u));
label.lft(btex 1 etex, (-1.5u, 0));


draw (-1.5u, 2u){dir 45}..(0, 1.8u)..{dir -45}(1.5u, 2u) withcolor green;
draw (-1.5u, 1u){dir 45}..{dir 45}(1.5u, 1u) withcolor red;
draw (-1.5u, 0){dir 45}..{dir -45}(1.5u, 0);

\stopMPpage
\stopbuffer

\placefigure[left][modell_atomhuelle:2025:01:29:skizze]{Modell}{
\typesetbuffer[modell_atomhuelle:2025:01:29:skizze]
}

Mögliche Wellenlängen


% \sidenote{}

\sidenote{
$n$ ist eine Zählvariable, Anzahl der Knoten ist gleich $n+1$
}
\startplaceformula[2025:01:29:wellenlaengen]
\startformula
\lambda_n = \frac{2L}{n}
\stopformula
\stopplaceformula

de Brogli:

\startplaceformula[2025:01:29:debrogli]
\startformula
p_n = \frac{h}{\lambda_n}=\frac{h\cdot n}{2L}
\stopformula
\stopplaceformula

\startformula
E_n=\frac{p_n^2}{2m}=\frac{h^2\cdot n^2}{8m \cdot L^2}=\frac{h^2}{8m}\cdot \frac{n^2}{L^2}
\stopformula

Daraus kann man die Wellenfunktion $\phi_n$ ermitteln.

Bedingung: 

\startplaceformula[2025:01:29:phi]
\startformula
\psi(x)+a\cdot \ddot{\psi}(x)=0
\stopformula
\stopplaceformula

mit: $a \in \reals$

Schrödingergleichung:

\startplaceformula[2025:01:29:schroedinger]
\startformula
\ddot{\psi}(x)=-\frac{8\cdot \pi^2 \cdot m}{h^2} \cdot (E - E_{\text{pot}}(x)) \cdot \psi(x)
\stopformula
\stopplaceformula

erfüllt mit $\sin-$ / $\cos-$ Fkt 

\startframedwarning
Abiturrelevant:
\stopframedwarning

$|\psi(x)|^2$ gibt die Aufenthaltswahrscheinlichkeit des $e^-$ am Ort $x$ an.

\startitemize[packed]
\item $n$: Quantenzahl -> Energieniveau
\item Bez. S. 235
\stopitemize

\startframedtipp
{\bf Orbital: S. 238 239}

Das "richtige Potenzial ist angegeben links in der abbildung, sog collomb
potenzial. Das beschreibt, wie die im endeffekt die anziehungskraft zwishen
zwei Ladungen im 3D Raum.  $n=1$ ist der Grundzustand. Der Unterschied zu dem and er Tafel (Potenzialtopf):

\startitemize[packed,m]
\item 3D: relativ simple
\item Hier angegebenes Potenzial realistischer
\item Anziehungskraft wird schwächer, desto weiter weg, niemals null
\item 
\stopitemize

Wenn wir n auf 2 setzen ist der abstand groeser
\stopframedtipp



\startframedtipp
DeBrogli muss ich kennen
Unter Quantenobjekte

und Herleitung
\stopframedtipp
